\section{本章小结}
\label{sec:meas-summary}

本章开头问：区间里的有理点有多长，挖去无数个中间段后剩下的点集有多长，怎样把“分开测量、合起来相加”推广到这些集合。计数、长度与体积共有的是有限可加性，Jordan测量用有限个矩形从内外夹逼，能处理边界可以用有限个小矩形盖住的集合。有理点集的内外测量分别为$0$与$1$，暴露了有限覆盖的局限；改用可数覆盖，总长度可以任意小。

由此得到两项要求：允许测量的集合族对可数运算封闭，测量规则对可数不交分割可加。$\sigma$代数规定测量范围，生成$\sigma$代数只补入封闭性迫使加入的集合，度量空间的开集生成Borel集，实轴上的区间或半直线已经足够生成它。测度在$\sigma$代数上非负、可数可加，由此推出单调性、可数次可加性与从下连续性；从上连续性需要一个测度有限的起点。零测集与几乎处处说明哪些例外可以忽略，完备化补入零测集的子集而不改变原有集合的测度。

测度的构造分三步。外测度用可数覆盖的最小费用给每个子集赋值，但一般不可加；Carathéodory判据选出与外测度相容的集合，它们构成$\sigma$代数，外测度在其上可数可加。从集合代数上的预测度出发，这一路线给出延拓；延拓的唯一性另需$\sigma$有限条件，并由$\pi$--$\lambda$定理在生成族上比较。

对单调不减且右连续的$F$，区间增量$F(b)-F(a)$确定唯一的Stieltjes测度，跳跃给出点质量；取$F(x)=x$得到Lebesgue测度，它平移不变，其完备化正是Lebesgue可测集。Cantor集零测而不可数，说明测度与基数是两回事；不可测集的存在说明测度只能定义在某个$\sigma$代数上。

有符号测度允许正负值，但至多取一种无穷，以免出现无穷减无穷。Hahn分解把空间切成正集与负集，切法不唯一；Jordan分解把有符号测度写成两个分别集中在不交集合上的测度之差，这个表示唯一且最小。全变差$\lvert\eta\rvert=\eta^++\eta^-$累计正负两方面的质量，也等于有限分割上绝对值之和的上确界，它的零测集正是有符号测度完全看不见的集合。

可测映射用原像规定哪些函数可以用测度研究，只需在生成族上检查；实值函数的可测性归结为阈值集合可测。连续函数都Borel可测，反之不然。可测函数对代数运算和可数次上下确界、上下极限封闭，扩展实值的运算要先说明有定义的范围；零测修改是否保持可测性取决于完备性。

按取值分层得到的递增简单函数列从下方逼近每个非负可测函数，这是第~\ref{chap:integration}章定义积分的出发点。两个测度怎样比较、何时能写成参考测度的加权积分，留到第~\ref{chap:measure-densities}章；质量沿可测映射的转移留到第~\ref{chap:pushforward-measures}章。
